\chapter*{LEMBAR PENGESAHAN}
\addcontentsline{toc}{chapter}{LEMBAR PENGESAHAN}


\begin{center}
    \vspace{\stretch{1}}
  {\large\bfseries PERANCANGAN DAN IMPLEMENTASI WAREHOUSE MANAGEMENT CONTROLLER \& SISTEM LOKALISASI UNTUK PROTOTIPE SMART WAREHOUSE}\\
     \vspace{\stretch{1}}

  {\Large \textbf{Laporan Tugas Akhir}}\\


  \vspace{\stretch{1}}


  {\large Oleh}\\[0.3cm]
    \textbf{
    {\large Raizan Iqbal Resi}\\
    {\large 18222068}
  }\\

  \vspace{0.5cm}

  {\large Program Studi Sistem dan Teknologi Informasi}\\
  {\large Sekolah Teknik Elektro dan Informatika}\\
  {\large Institut Teknologi Bandung}\\

  \vspace{\stretch{1}}

  Laporan Tugas Akhir ini telah disetujui dan disahkan\\
  di Bandung, pada tanggal \today\\[1cm]

\begin{tabular}{p{1cm}p{7cm}p{7cm}}
  & Pembimbing 1 & Pembimbing 2 \\[2cm]
  & Dr.techn. Ary Setijadi Prihatmanto, S.T., M.T. & I Gusti Bagus Baskara Nugraha, S.T., M.T., Ph.D. \\[0.2cm]
  & NIP. 197208271997021003 & NIP. 9197601242010121001
\end{tabular}

\vspace{\stretch{1}}
\end{center}
\noindent

% ==========================================
% Jika ada 2 pembimbing TA, uncomment dan edit 
% tabular di bawah ini. Kemudian, comment out atau hapus
% bagian versi 1 pembimbing di atas.
% ==========================================

